\uuid{CkuL}
\exo7id{3471}
\titre{exo7 3471}
\auteur{quercia}
\organisation{exo7}
\datecreate{2010-03-10}
\isIndication{false}
\isCorrection{false}
\chapitre{Déterminant, système linéaire}
\sousChapitre{Système linéaire, rang}
\cpge{mpsi-mp2i.matrices-applications-lineaires.b, pcsi.matrices-applications-lineaires.b, ptsi.matrices-applications-lineaires.b, mpsi-mp2i.calcul-matriciel-systemes.a, pcsi.calcul-matriciel-systemes-lineaires.a, ptsi.calcul-matriciel-systemes-lineaires.a, bcpst1.matrices, bcpst1.applications-lineaires-matrices, pc.complements-espaces-vectoriels.b, psi.complements-espaces-vectoriels.b, pt.complements-algebre-lineaire.b, bcpst2.applications-lineaires-matrices.c}
\module{Algèbre}
\niveau{L2}
\difficulte{}

\contenu{
\question{
Soient $A \in \mathcal{M}_{n,p}(K)$ et $B \in \mathcal{M}_{n,q}(K)$. On considère
$C = (A\ B) \in \mathcal{M}_{n,p+q}(K)$.

Montrer que : $\mathrm{rg}(C) = \mathrm{rg}(A) \Leftrightarrow \exists\ P \in \mathcal{M}_{p,q}(K)$ tq $B = AP$.
}
}
